% ============================================================================
% AD7031 Homework 2 -- covers Weeks 4, 6, 8 (mean and covariance of a random
% vector, LP duality, conic duality, risk measures VaR and CVaR).
% 100 points, 4 problems of 25.
% Self-contained on purpose -- no \input{preamble} -- compiles with pdflatex.
% ============================================================================
\documentclass[10pt]{article}

\usepackage[a4paper, margin=2.2cm]{geometry}
\usepackage{amsmath, amssymb}
\usepackage[dvipsnames]{xcolor}
\usepackage{array}
\usepackage{enumitem}
\usepackage[hidelinks]{hyperref}

\definecolor{ruleink}{RGB}{18, 68, 140}
\definecolor{keyink}{RGB}{20, 90, 20}

\setlength{\parindent}{0pt}
\pagestyle{plain}

\AtBeginDocument{%
  \abovedisplayskip=5pt plus 1pt
  \belowdisplayskip=5pt plus 1pt
  \abovedisplayshortskip=3pt
  \belowdisplayshortskip=3pt}

% ---- one problem -----------------------------------------------------------
\newcommand{\problem}[2]{%   points, title
  \par\bigskip
  {\color{ruleink}\rule{\textwidth}{1pt}}\par\nobreak\vspace{3pt}
  {\large\bfseries #2}\hfill{\normalsize\bfseries(#1 points)}\par\nobreak\vspace{4pt}}

% ---- one-line hint, blue italic --------------------------------------------
\newcommand{\hint}[1]{%
  \par\vspace{3pt}{\color{ruleink}\emph{hint:\ #1}}\par}

\newlist{parts}{enumerate}{1}
\setlist[parts]{label=(\alph*), leftmargin=2.2em, itemsep=5pt}

\begin{document}

\begin{center}
  {\large AD7031 Optimization Under Uncertainty}\\[3pt]
  {\Large\bfseries Homework 2}\\[4pt]
  {\small \textbf{Due Oct 31, start of class}}
\end{center}
\medskip\hrule\medskip

Write down your solutions and submit them as a single PDF to
\texttt{lossof0706@gmail.com}.\\
\textbf{\underline{Your solutions must be handwritten.}}

% ================================================================ Problem 1
\problem{25}{Problem 1. Mean and covariance of a random vector}

Three months of returns are observed, each equally likely ($p_s = \tfrac13$):

\begin{center}
\begin{tabular}{c|ccc}
           & Tech   & Energy & Bond   \\ \hline
 $\xi^1$   & $0.00$ & $0.04$ & $0.03$ \\
 $\xi^2$   & $0.06$ & $0.10$ & $0.03$ \\
 $\xi^3$   & $0.12$ & $0.07$ & $0.03$ \\
\end{tabular}
\end{center}

\begin{parts}
\item \textbf{(10 pts)} \textbf{By hand}, compute the mean $\mu = \mathbb{E}[\tilde\xi]$ and the
  covariance matrix $\Sigma = \mathrm{Cov}(\tilde\xi)$, and state the dimensions of $\mu$ and $\Sigma$.

\item \textbf{(10 pts)} Suppose your optimal portfolio is $x^{\star} = (x_1^{\star}, x_2^{\star}, x_3^{\star})^{\top}$.
  From the table, compute, in terms of $x^{\star}$,
  \begin{itemize}[leftmargin=2.2em, itemsep=2pt, topsep=3pt]
    \item the mean $\mathbb{E}[\tilde\xi^{\top}x^{\star}]$;
    \item the variance $\mathrm{Var}[\tilde\xi^{\top}x^{\star}]$.
  \end{itemize}

\item \textbf{(5 pts)} Show that
  \begin{itemize}[leftmargin=2.2em, itemsep=2pt, topsep=3pt]
    \item $\mathbb{E}[\tilde\xi^{\top}x^{\star}] = \mu^{\top}x^{\star}$;
    \item $\mathrm{Var}[\tilde\xi^{\top}x^{\star}] = x^{\star\top}\Sigma\, x^{\star}$.
  \end{itemize}
  \hint{we discussed this in class.}
\end{parts}

% ================================================================ Problem 2
\problem{25}{Problem 2. Linear program duality}

Consider the following primal linear program and its dual, in the canonical form of the lecture note:
\[
\text{(P)} \quad
\begin{aligned}
  \min_{x \in \mathbb{R}^n} \quad & c^{\top}x \\
  \text{s.t.} \quad & Ax \le b
\end{aligned}
\qquad\qquad
\text{(D)} \quad
\begin{aligned}
  \max_{\theta \in \mathbb{R}^m} \quad & -b^{\top}\theta \\
  \text{s.t.} \quad & c + A^{\top}\theta = 0 \\
                    & \theta \ge 0
\end{aligned}
\]

\begin{parts}
\item \textbf{(8 pts)} Derive the dual problem (D) from the primal problem (P). \textbf{Show all the steps clearly.}

\item \textbf{(7 pts)} Consider the primal problem
\[
\begin{aligned}
  \min_{x \in \mathbb{R}^2} \quad & -4x_1 - 5x_2 \\
  \text{s.t.} \quad & 3x_1 + 2x_2 \le 60 \\
                    & x_1 + 3x_2 \le 48 \\
                    & x_1 \ge 3 \\
                    & x_1,\ x_2 \ge 0 .
\end{aligned}
\]
  \begin{itemize}[leftmargin=2.2em, itemsep=2pt, topsep=3pt]
    \item Give $c$, $A$, $b$.
    \item Write the dual problem (D) numerically.
  \end{itemize}

\item \textbf{(7 pts)} Derive the primal problem (P) from the dual problem (D). \textbf{Show all the steps clearly.}
  \hint{nothing really changes from the derivation in (a).}

\item \textbf{(3 pts)} Let $x^{\star}$ and $\theta^{\star}$ be optimal solutions of the primal problem (P) and the dual problem (D),
  respectively. What can you say about $c^{\top}x^{\star}$ and $b^{\top}\theta^{\star}$?
\end{parts}

% ================================================================ Problem 3
\problem{25}{Problem 3. Conic program duality}

Consider
\[
\begin{aligned}
  \max_{x \in \mathbb{R}^n} \quad & \mu^{\top}x \\
  \text{s.t.} \quad & x^{\top}\Sigma x \le \sigma^2 \\
                    & e^{\top}x \le 1 \\
                    & x \ge 0
\end{aligned}
\]
where $e = [1, 1, \dots, 1]^{\top}$ is the all-ones vector and $\sigma$ is the risk cap.

\begin{parts}
\item \textbf{(10 pts)} Recall the second-order cone
  \[
    \mathrm{SOC}_Q^{n+1} = \big\{ (x, t) \in \mathbb{R}^n \times \mathbb{R} \;:\; \sqrt{x^{\top}Q\,x} \le t \big\},
    \qquad Q \succ 0 .
  \]
  Write the model in the conic form $\min c^{\top}z$ s.t.\ $Az \le b$, $z \in C$.
  \begin{itemize}[leftmargin=2.2em, itemsep=2pt, topsep=3pt]
    \item Give $z$ (the decision variable), $c$, $A$, $b$.
    \item Write $C$.
  \end{itemize}
  \hint{use $z = [x,\ t]^{\top} \in \mathbb{R}^{n+1}$, then write out $c^{\top}z$ and $Az \le b$ to check that you recover the original problem.}

\item \textbf{(10 pts)} With $\Sigma = I$, the dual problem of the second-order cone program in (a) is
  \[
  \begin{aligned}
    \max \quad & -\theta_1 - \sigma\, \theta_2 \\
    \text{s.t.} \quad & \theta_1,\ \theta_2 \ge 0, \quad \theta_3 \in \mathbb{R}^n_+ \\
                      & \|\theta_1 e - \mu - \theta_3\|_2 \;\le\; \theta_2 .
  \end{aligned}
  \]
  Derive it. \textbf{Show all the steps clearly.}

\item \textbf{(5 pts)} State a sufficient condition for strong duality. For which values of
  $\sigma$ does it hold for the primal problem in (a)?
\end{parts}

% ================================================================ Problem 4
\problem{25}{Problem 4. Risk measures}

Recall the general stochastic program
\[
\text{(SP)} \quad
\begin{aligned}
  \min_{x \in \mathbb{R}^n} \quad & \mathcal{R}_{\mathbb{P}}\big[f_0(x, \tilde\xi)\big] \\
  \text{s.t.} \quad & \mathcal{R}_{\mathbb{P}}\big[f_i(x, \tilde\xi)\big] \le 0 \qquad \forall i \in [I],
\end{aligned}
\]
where $\tilde\xi \sim \mathbb{P}$ and $\mathcal{R}_{\mathbb{P}}$ is a risk measure of our choice.

\begin{parts}
\item \textbf{(5 pts)} Suppose we choose $\mathcal{R}_{\mathbb{P}} = \mathrm{VaR}_{\varepsilon}$, i.e., the objective function in (SP) is
  \[
  \begin{aligned}
    \mathrm{VaR}_{\varepsilon}\big[f_0(x, \tilde\xi)\big] \;=\; \min_{t \in \mathbb{R}} \quad & t \\
    \text{s.t.} \quad & \mathbb{P}\big(f_0(x, \tilde\xi) \le t\big) \ge 1 - \varepsilon .
  \end{aligned}
  \]
  Using this minimization problem, explain in plain Korean what $\mathrm{VaR}_{\varepsilon}\big[f_0(x, \tilde\xi)\big]$ means.

\item \textbf{(5 pts)} Show the reformulation of the constraint $\mathrm{VaR}_{\varepsilon}\big[f_i(x, \tilde\xi)\big] \le 0$.
  \hint{use
    \[
      \mathrm{VaR}_{\varepsilon}\big[f_i(x, \tilde\xi)\big] \le 0
      \quad \Longleftrightarrow \quad
      \left[\begin{aligned}
        \min_{t_i \in \mathbb{R}} \quad & t_i \\
        \textnormal{s.t.} \quad & \mathbb{P}\big(f_i(x, \tilde\xi) \le t_i\big) \ge 1 - \varepsilon
      \end{aligned}\right] \le 0
      \quad \Longleftrightarrow \quad
      \begin{aligned}
        & t_i \le 0 \\
        & \mathbb{P}\big(f_i(x, \tilde\xi) \le t_i\big) \ge 1 - \varepsilon .
      \end{aligned}
    \]}

\item \textbf{(5 pts)} Suppose we choose $\mathcal{R}_{\mathbb{P}} = \mathrm{CVaR}_{\varepsilon}$, i.e., the objective function in (SP) is
  \[
    \mathrm{CVaR}_{\varepsilon}\big[f_0(x, \tilde\xi)\big] \;=\; \min_{\beta \in \mathbb{R}}\
      \beta + \frac{1}{\varepsilon}\, \mathbb{E}_{\mathbb{P}}\big[\max\{f_0(x, \tilde\xi) - \beta,\ 0\}\big] .
  \]
  Using this minimization problem, explain in plain Korean what $\mathrm{CVaR}_{\varepsilon}\big[f_0(x, \tilde\xi)\big]$ means.

\item \textbf{(5 pts)} Write the reformulation of (SP) with $\mathcal{R}_{\mathbb{P}} = \mathrm{CVaR}_{\varepsilon}$.
  \textbf{Clearly show the reformulation of the objective
  $\mathrm{CVaR}_{\varepsilon}\big[f_0(x, \tilde\xi)\big]$ and of every constraint
  $\mathrm{CVaR}_{\varepsilon}\big[f_i(x, \tilde\xi)\big] \le 0$, $i \in [I]$.}
  \hint{do the same as in (b).}

\item \textbf{(5 pts)} For computational tractability, we may prefer CVaR over VaR. Explain why.
\end{parts}

\end{document}
